\ExamPart{Trắc nghiệm nhiều phương án lựa chọn}{3.0}{Thí sinh trả lời từ câu 1 đến câu 12. Mỗi câu thí sinh chỉ chọn một phương án.}

\NQuestion Tam thức bậc hai $f\left(x\right)=x^2-4x+3$ âm khi
\ChoicesOneLine
{$x\in(-\infty;1)\cup(3;+\infty)$.}
{$x\in(1;3)$.}
{$x\in[1;3]$.}
{$x\in(-\infty;1]\cup[3;+\infty)$.}
{2}
\SolutionBlock{Ta có $x^2-4x+3=(x-1)(x-3)$. \\
    Vì hệ số $a=1>0$ nên tam thức bậc hai âm giữa hai nghiệm phân biệt $1$ và $3$. \\
    Do đó $f\left(x\right)<0$ khi $x\in(1;3)$.
}

\NQuestion Tập nghiệm của bất phương trình $(x-2)(x+1)\le 0$ là
\ChoicesOneLine
{$(-\infty;-1]\cup[2;+\infty)$.}
{$[-1;2]$.}
{$(-1;2)$.}
{$(-\infty;-1)\cup(2;+\infty)$.}
{2}
\SolutionBlock{Hai nghiệm của vế trái là $x=-1$ và $x=2$. \\
    Vì tích hai thừa số bậc nhất có hệ số đầu dương nên biểu thức không dương trên đoạn nằm giữa hai nghiệm, kể cả hai đầu mút. \\
    Vậy tập nghiệm là $[-1;2]$.}

\NQuestion Phương trình $x^4-5x^2+4=0$ có số nghiệm thực là
\ChoicesOneLine{$1$.}{$2$.}{$3$.}{$4$.}{4}
\SolutionBlock{Đặt $t=x^2\ (t\ge 0)$. \\
    Khi đó ta được phương trình $t^2-5t+4=0\Leftrightarrow (t-1)(t-4)=0$. \\
    Suy ra $t=1$ hoặc $t=4$, nên $x=\pm1,\pm2$. \\
    Có tất cả $4$ nghiệm thực.}

\NQuestion Một bạn học sinh có $3$ áo, $2$ quần và $2$ đôi giày. Số cách chọn một bộ gồm $1$ áo, $1$ quần và $1$ đôi giày là
\ChoicesOneLine{$7$.}{$10$.}{$12$.}{$14$.}{3}
\SolutionBlock{Áp dụng quy tắc nhân: có $3$ cách chọn áo, $2$ cách chọn quần và $2$ cách chọn giày. \\
    Vậy số cách là $3\cdot2\cdot2=12$.}

\NQuestion Số cách xếp $5$ học sinh phân biệt ngồi vào một hàng ghế là
\ChoicesOneLine{$25$.}{$60$.}{$100$.}{$120$.}{4}
\SolutionBlock{Sắp xếp $5$ học sinh phân biệt vào một hàng là số hoán vị của $5$ phần tử: $P_5=5!=120$.}

\NQuestion Từ các chữ số $1,2,3,4,5$, lập được bao nhiêu số tự nhiên có $3$ chữ số khác nhau và chia hết cho $5$?
\ChoicesOneLine{$6$.}{$12$.}{$18$.}{$24$.}{2}
\SolutionBlock{Muốn chia hết cho $5$ thì chữ số hàng đơn vị phải là $5$. \\
    Chọn chữ số hàng trăm từ $\{1,2,3,4\}$ có $4$ cách. \\
    Chọn chữ số hàng chục từ $3$ chữ số còn lại có $3$ cách. \\
    Do đó có $4\cdot3=12$ số.
}

\NQuestion Trong mặt phẳng tọa độ $Oxy$, cho $A(-2;3),\ B(4;-1)$. Tọa độ vectơ $\overrightarrow{AB}$ là
\ChoicesOneLine
{$(2;-4)$.}
{$(6;-4)$.}
{$(-6;4)$.}
{$(6;4)$.}
{2}
\SolutionBlock{Ta có $\overrightarrow{AB}=(4-(-2);-1-3)=(6;-4)$.}

\NQuestion Trong mặt phẳng tọa độ $Oxy$, trung điểm của đoạn thẳng nối hai điểm $A(1;2)$ và $B(5;-4)$ là
\ChoicesOneLine{$(3;-1)$.}{$(2;-1)$.}{$(3;1)$.}{$(6;-2)$.}{1}
\SolutionBlock{Gọi $I$ là trung điểm của $AB$. Khi đó
\[
I\left(\dfrac{1+5}{2};\dfrac{2+(-4)}{2}\right)=(3;-1).
\]}

\NQuestion Trong mặt phẳng tọa độ $Oxy$, cho tam giác $ABC$ với $A(1;2),\ B(4;-1),\ C(-2;3)$. Tọa độ trọng tâm $G$ của tam giác là
\ChoicesOneLine
{$\left(1;\dfrac{4}{3}\right)$.}
{$\left(1;\dfrac{2}{3}\right)$.}
{$\left(\dfrac{4}{3};1\right)$.}
{$\left(\dfrac{1}{3};4\right)$.}
{1}
\SolutionBlock{Trọng tâm tam giác có tọa độ
\[
G\left(\dfrac{1+4+(-2)}{3};\dfrac{2+(-1)+3}{3}\right)=\left(1;\dfrac{4}{3}\right).
\]}

\NQuestion Đường thẳng đi qua $M(1;-2)$ và song song với đường thẳng $d:2x-y+3=0$ có phương trình là
\ChoicesOneLine
{$2x-y+1=0$.}
{$2x-y-4=0$.}
{$x+2y+3=0$.}
{$2x+y-4=0$.}
{2}
\SolutionBlock{Đường thẳng cần tìm song song với $d$ nên có dạng $2x-y+c=0$. Thay điểm $M(1;-2)$ vào được $2\cdot1-(-2)+c=0\Rightarrow c=-4$. Vậy phương trình là $2x-y-4=0$.}

\NQuestion Khoảng cách từ điểm $P(1;2)$ đến đường thẳng $d:3x+4y-12=0$ bằng
\ChoicesOneLine{$1$.}{$\dfrac{1}{5}$.}{$\dfrac{2}{5}$.}{$\dfrac{3}{5}$.}{2}
\SolutionBlock{Khoảng cách từ điểm $P(x_0;y_0)$ đến đường thẳng $ax+by+c=0$ là
\[
d(P,d)=\dfrac{|ax_0+by_0+c|}{\sqrt{a^2+b^2}}.
\]
Suy ra
\[
d(P,d)=\dfrac{|3\cdot1+4\cdot2-12|}{\sqrt{3^2+4^2}}=\dfrac{1}{5}.
\]}

\NQuestion Gọi $\alpha$ là góc giữa hai đường thẳng $d_1:x-y+1=0$ và $d_2:x+2y-4=0$. Khi đó
\ChoicesOneLine
{$\tan\alpha=\dfrac{1}{3}$.}
{$\tan\alpha=1$.}
{$\tan\alpha=3$.}
{$\tan\alpha=\sqrt3$.}
{3}
\SolutionBlock{Ta có $d_1:x-y+1=0\Rightarrow y=x+1$ nên hệ số góc $m_1=1$. Đồng thời $d_2:x+2y-4=0\Rightarrow y=-\dfrac12x+2$ nên $m_2=-\dfrac12$. Khi đó
\[
\tan\alpha=\left|\dfrac{m_2-m_1}{1+m_1m_2}\right|=\left|\dfrac{-\dfrac12-1}{1-\dfrac12}\right|=3.
\]}
